\chapter{绪论}
\label{cha:intro}

\section{研究背景与意义}

\noindent
脑机接口能够在不依赖外周神经和肌肉输出的条件下建立人脑与外部设备之间的信息通道，在神经康复、辅助交互、智能可穿戴和新型人机协同等方向具有重要应用价值 \cite{lotte2018review}。其中，基于脑电（Electroencephalography, EEG）的非侵入式脑机接口因采集成本较低、使用风险较小、便于重复实验和系统集成，成为当前研究与应用开发最活跃的技术路线之一。运动想象（Motor Imagery, MI）任务由于无需真实肢体动作即可诱发可识别的感觉运动节律变化，长期被视为脑机接口研究中的典型范式。

\noindent
然而，面向真实使用场景的脑机接口系统与标准公开数据集上的离线解码任务并不完全等价。便携式智能终端、头戴式交互设备和低成本辅助系统通常要求前端具备少导联、轻量化、低功耗和无线通信等特征；相应地，后端解码模型也需要在有限算力和有限存储条件下完成稳定推理。对于这类场景而言，单纯追求高密度导联、高复杂度模型和受控协议下的更高离线分数，并不必然能够转化为真实系统中的可部署性与可用性。

\noindent
近几年，脑电解码研究一方面沿着更深的卷积网络、更复杂的自注意力结构以及更大规模训练数据持续推进，另一方面也逐步显露出与真实低通道场景之间的错位。高密度标准条件下得到的优势模型，在导联数显著下降、个体差异增强、试次波动增大和端侧资源受限的条件下，往往难以直接复用。对于面向双导、低成本设备的脑机接口原型而言，更关键的问题并不是继续追求“更大模型”，而是如何在真实约束下找到合理的任务定义、够用的模型结构以及可落地的系统链路。

\noindent
基于上述问题意识，本文围绕双导、便携、低成本端侧脑机接口的真实需求，研究运动想象脑电解码任务如何从高密度标准离线场景收束为可部署的低通道系统方案。论文一方面完成双导脑电采集硬件系统与通信链路实现，另一方面围绕低通道端侧场景设计改进型自注意力解码方法，并进一步通过从 `22导 / 4分类` 到 `2导 / 2分类` 的问题收束、离线验证与在线联动初步验证，形成从问题定义到系统落地的完整研究路径。

\section{相关研究现状}
\subsection{高密度离线条件下的脑电解码研究现状}

在标准化运动想象脑电解码任务中，传统机器学习方法构成了长期的重要基线。以共空间模式（Common Spatial Pattern, CSP）及其滤波器组扩展方法（Filter Bank CSP, FBCSP）为代表的空间滤波路线，能够突出不同类别在空间方差分布上的差异，并与线性判别分析（LDA）、支持向量机（SVM）等浅层分类器结合取得较高的分类效率 \cite{ang2012filter, lotte2018review}。这类方法在标准数据集和受控实验条件下具有实现简单、可解释性较强等优点，但其效果往往较为依赖预设频带、人工特征和稳定采集条件。

\noindent
为减少对手工特征工程的依赖，卷积神经网络随后成为脑电解码的重要技术路线。Schirrmeister 等提出的深浅卷积网络说明，端到端学习方式可以直接从原始脑电信号中提取具有判别性的时空特征 \cite{schirrmeister2017deep}；Lawhern 等提出的 EEGNet 则以紧凑型卷积结构兼顾了较低参数规模与较好的跨范式适应能力 \cite{lawhern2018eegnet}。相比传统特征工程流程，CNN 有效推动了脑电解码向自动化表示学习演进。

\noindent
在 CNN 的基础上，Transformer 及其变体被引入脑电解码领域，以增强对长程时序依赖和跨特征关联的建模能力 \cite{vaswani2017attention}。Song 等提出的 EEG Conformer 通过卷积与自注意力协同提取局部时空特征和全局时间依赖 \cite{songEEG2022}，Xie 等进一步展示了 Transformer 结构在原始脑电分类任务中的应用潜力 \cite{xie_transformer-based_2022}。随后，面向脑电序列的高效 Transformer 设计逐步出现，研究重点从“能否使用自注意力”转向“如何在控制复杂度的同时保留关键建模能力”。例如，Wan 等将通道注意力与窗口化机制结合建模时空耦合信息 \cite{wan2023novel}，Hua 等通过卷积与滑动窗口注意力的协同设计增强了多尺度特征建模能力 \cite{hua2023improved}，Keutayeva 等则进一步强调了紧凑型结构在跨被试任务中的部署潜力 \cite{keu2024compact}。

\noindent
综合来看，高密度离线条件下的脑电解码研究已经形成了从传统空间滤波、卷积网络到自注意力建模的清晰演进路径，并在公开数据集上不断提升分类性能。然而，这一研究脉络主要建立在较多导联、标准化协议和离线评估框架之上，其优越性是否能够在低通道、低成本和端侧场景下保持，仍需要进一步验证。

\subsection{低通道与端侧部署研究现状}

随着可穿戴设备和便携式脑机接口的发展，低通道解码和端侧部署逐渐成为新的研究重点。一类工作从采集端出发，尝试通过减少导联数、优化电极布局和简化采集链路来提升系统便携性；另一类工作则从模型端出发，围绕参数压缩、量化推理和低功耗实现展开探索。在运动想象脑电场景中，已有研究开始关注紧凑型模型在边缘设备上的运行可行性。例如，Wan 等、Schiavone 等和 Ene 等分别从低功耗边缘计算、8 bit 量化实现和 Arduino Nano 级平台落地等角度，验证了紧凑型 MI 解码模型在资源受限平台上的应用潜力 \cite{Wan20, Sch20, Wan22b, Ene23}。这些工作说明，面向端侧的 MI 脑电解码并非不可实现，但现阶段的直接证据仍主要集中在 CNN 或更简单结构上。

\noindent
与此同时，紧凑型 Transformer 和相关时序模型在端侧平台上的部署证据也开始积累，但其分布并不均衡。相邻任务中的研究已经给出了 MCU 平台上 Transformer 类模型的真实延迟与能耗数据 \cite{bus2024reducing}，部分嵌入式时序任务也开始面向 RK3568、NPU 或异构推理引擎进行优化 \cite{You24, Ism23}。不过，对于本文更关心的“运动想象脑电 + 低通道 + 自注意力模型 + RK3568/RKNN 原型节点”这一组合，公开文献中仍缺少足够完整的端到端实测结果。现有工作更多回答了“模型能否被压缩”或“平台是否具有推理能力”，而较少同时回答“在真实低通道脑电信号条件下，什么模型与什么任务定义才真正适合部署”。

\noindent
此外，低通道和端侧场景并不仅仅意味着导联数更少或模型更小，还意味着系统级约束发生变化。前端采集噪声、半电池效应引起的慢漂移、BLE 传输的时序波动、被试个体差异以及在线校准成本，都会共同影响最终的解码可用性。因此，低通道端侧脑机接口研究需要从“单点优化”走向“任务定义、模型设计与系统实现的协同约束”，这也是本文后续工作的直接出发点。

\section{现有研究存在的不足}

综合前述研究可以看出，当前运动想象脑电解码虽然在模型设计、局部高效注意力和轻量化探索方面取得了明显进展，但对于双导、低成本和端侧约束下的真实系统目标而言，仍存在以下三方面不足。

\noindent
第一，现有高密度标准任务与低通道真实场景之间存在明显错位。公开数据集中的高密度四分类设置为算法研究提供了统一基准，但当导联数显著降低时，可观测信息边界随之收缩，原有任务定义并不一定仍然合理。许多工作默认“模型只要足够强就能迁移到更少通道”，却较少系统回答在双导条件下究竟应保留什么目标、舍弃什么目标，以及应如何构建相应的定量证据链。

\noindent
第二，算法设计与端侧部署约束之间的协同仍然不足。现有研究往往分别强调高离线精度、模型压缩或硬件加速中的某一方面，但较少在统一框架下同时考虑低通道输入、结构表达能力、参数规模、量化友好性和推理延迟等约束 \cite{keu2024compact, wan2023novel, hua2023improved}。尤其是在运动想象脑电场景中，“能在公开数据集上取得较高准确率”并不等同于“能够在双导端侧系统中稳定运行并保持可用表现”。

\noindent
第三，硬件实现、在线联动与解码模型之间仍缺少完整闭环。便携式脑机接口不仅取决于模型推理性能，还同时受到采集链路噪声、BLE 传输时序、试次间漂移和在线校准代价等因素影响。现有工作中，真实在线 MI 证据总体仍然有限，伪在线评估与真实系统试跑之间也常常缺乏清晰边界。如何在双导硬件约束下，将任务收束、模型设计、前端采集和在线联动机制统一起来，是本文试图解决的核心问题。

\section{本文主要研究内容与技术路线}
\subsection{主要研究内容}

围绕上述问题，本文面向双导、便携、低成本端侧脑机接口场景，开展了以下几方面研究工作。

\noindent
（1）设计并实现双导脑电采集硬件系统。针对低成本、便携式采集需求，本文完成了以双导联采集为核心的硬件系统设计与实现，构建了脑电前端、主控、BLE 回传和上位机交互的基础链路，并通过节律验证、标签链路验证和通信稳定性测试，为后续在线实验提供系统基础。

\noindent
（2）研究面向低通道端侧场景的改进型自注意力脑电解码方法。针对全局自注意力复杂度较高、难以直接适配端侧平台的问题，本文设计了兼顾局部时序建模能力与部署友好性的改进型自注意力解码方法，在标准条件下完成性能、复杂度与量化友好性分析，并为后续低通道扩展提供算法基础。

\noindent
（3）构建从高密度标准任务到双导可落地任务的收束路径。围绕 `22导 / 4分类` 标准设置，本文系统分析了导联压缩下的性能退化边界，并在此基础上将目标进一步收束到 `2导 / 2分类` 的左右手运动想象解码任务，从而使算法目标与双导系统可观测信息边界相匹配。

\noindent
（4）完成双导端侧系统集成与在线联动初步验证。本文以双导硬件为前端、以上位机和可扩展推理节点为原型载体，构建了预训练冷启动、本地数据累计、轮次训练与下一轮推理呈现的在线联动机制，并结合伪在线评估和真实在线试跑，对系统的可实现性进行了初步验证。需要指出的是，当前实验阶段的在线验证载体主要为 Linux PC 原型环境，后续目标部署节点则面向 RK3568 / RKNN 一类嵌入式平台。

\subsection{技术路线与论文组织}

全文的技术路线与组织结构如图\ref{fig:org}所示。
\begin{figure}[htbp]
  \centering
  % \includegraphics[width=0.8\textwidth]{org_structure}
  \caption{论文技术路线与组织结构}
  \label{fig:org}
\end{figure}

\noindent
本文后续各章安排如下：

\noindent
第 2 章介绍运动想象脑电解码与端侧系统实现所需的关键技术基础，包括 MI 脑电生理特性、预处理、自注意力与高效注意力机制，以及端侧部署、BLE 与采集前端相关知识。

\noindent
第 3 章围绕双导脑电采集硬件系统的设计与实现展开，重点说明采集前端、主控、无线通信和基础链路验证结果。

\noindent
第 4 章研究面向低通道端侧场景的改进型自注意力脑电解码方法，在标准高密度条件下建立离线基线，并进一步分析导联压缩下的性能退化边界，完成从 `22导 / 4分类` 到 `2导 / 2分类` 的任务收束与模型适配。

\noindent
第 5 章进一步将第 4 章得到的低通道解码主线映射到双导系统中，介绍系统集成、在线实现机制、伪在线评估以及真实在线初步验证结果。

\noindent
第 6 章对全文工作进行总结，并围绕更强的人机协同、自适应学习和更自然的实验范式给出后续展望。
